\begin{abstract}
Automated leukocyte morphology classifiers are commonly evaluated under a closed-set assumption, although deployed or reused models may encounter morphologies outside the training taxonomy. We evaluated whether high closed-set performance implies reliable rejection of unseen morphologies under internal and external domain shift. The known taxonomy comprised basophils, eosinophils, lymphocytes, monocytes, and segmented neutrophils. In MLL23, 12,793 known images were used for training, 2,741 for validation, 2,742 for known testing, and 23,345 held-out images from 13 morphologies for unknown testing. ImageNet-pretrained ResNet18 models trained with cross entropy or ArcFace were assessed across five seeds and two split protocols using post-hoc open-set scores. Internal closed-set macro-F1 was high for both losses, approximately 0.97. Internal open-set recognition was lower: CE+ViM achieved AUROC 0.877 $\pm$ 0.004 on V1 and 0.874 $\pm$ 0.008 on the near-duplicate-aware V2 sensitivity split, with FPR@95TPR 0.484 $\pm$ 0.027 and 0.528 $\pm$ 0.050. On AML-Cytomorphology\_LMU, an independent 18,365-image test-only dataset, performance degraded substantially; external AUROC ranged from 0.563 to 0.696 on V1 systems and from 0.573 to 0.696 on V2 systems, with high FPR@95TPR. Failure was morphology-dependent, with neutrophil-band, monoblast, myeloblast, and atypical lymphocyte among the hardest external unknowns. ArcFace improved some known-class geometry summaries but did not reproduce as a robust superior open-set representation. Persistent-homology analyses were negative or explanatory rather than predictive. Strong closed-set leukocyte classification therefore did not ensure reliable unseen-morphology rejection, especially under morphological similarity and external domain shift.
\end{abstract}
